%%%%%%%%%%%%%%%%%%%%%%%%%%%%%%%%%%%%%%%%%%%%%%%%%%%%%%%%%%%%%%%%%%%%%%
% Evolução da Maturidade de Desenvolvimento em Projetos Spring Boot
% TCC - Engenharia de Software - PUC Minas
%
% Para compilar:
%   pdflatex Article.tex
%   bibtex   Article
%   pdflatex Article.tex
%   pdflatex Article.tex
%
% Observação: o arquivo de estilo `sbc-template.sty` deve estar no mesmo
% diretório (download em [https://www.sbc.org.br/documentos-da-sbc/](https://www.sbc.org.br/documentos-da-sbc/)
% summary/169-templates-para-artigos-e-capitulos-de-livros/878-modelosparapublicaodeartigos).
%%%%%%%%%%%%%%%%%%%%%%%%%%%%%%%%%%%%%%%%%%%%%%%%%%%%%%%%%%%%%%%%%%%%%%

\documentclass[12pt]{article}

\usepackage{sbc-template}
\usepackage{graphicx,url}
\usepackage{float}
\usepackage[brazil]{babel}
\usepackage{tabularx}
\newcolumntype{C}{>{\centering\arraybackslash}X}
\usepackage[utf8]{inputenc}
\usepackage[inline]{enumitem}
\usepackage{tcolorbox}
\usepackage{tikz}
\usetikzlibrary{arrows.meta,positioning}
\usepackage{amsmath}

%%%%%%%%%%%%%%%%% Pacotes para inserção de algoritmos - BEGIN
\usepackage[ruled,lined]{algorithm2e}

\usepackage{xcolor}
\definecolor{verde}{rgb}{0.25,0.5,0.35}
\definecolor{jpurple}{rgb}{0.5,0,0.35}
\definecolor{darkgreen}{rgb}{0.0, 0.2, 0.13}

\usepackage{listings}
\lstset{
    language=Java,
    basicstyle=\ttfamily\small,
    keywordstyle=\color{jpurple}\bfseries,
    stringstyle=\color{red},
    commentstyle=\color{verde},
    morecomment=[s][\color{blue}]{/**}{*/},
    extendedchars=true,
    showspaces=false,
    showstringspaces=false,
    numbers=left,
    numberstyle=\tiny,
    breaklines=true,
    backgroundcolor=\color{white!10},
    breakautoindent=true,
    captionpos=b,
    xleftmargin=0pt,
    tabsize=2
}
%%%%%%%%%%%%%%%%% Pacotes para inserção de algoritmos - END

\graphicspath{%
  {imgs/}%
  {Sections/imgs/}%
  {../Divulgacao/Apresentacao/imgs/}%
  {./}%
}

\sloppy

\title{Evolução da Maturidade de Desenvolvimento\\em Projetos Spring Boot}

\author{Gabriel Ferreira\inst{1}, Guilherme Augusto\inst{1}}

\address{Engenharia de Software --
  Pontifícia Universidade Católica de Minas Gerais (PUC-MG)\\
  Belo Horizonte -- MG -- Brazil
  \email{\{gabriel.amaral.1447804,gajsouza\}@sga.pucminas.br}
}

\begin{document}

\maketitle

\begin{abstract}
Internal software quality in object-oriented systems can be observed through structural metrics of coupling, complexity, inheritance, cohesion and size. In the Java ecosystem, Spring Boot is a widely adopted framework, yet longitudinal evidence on how the structural quality of public GitHub projects evolves---and whether the generative-AI era (2022--present) coincides with measurable change---remains scarce. This work aims to analyse that evolution across four temporal cohorts defined by repository creation date. Object-oriented metrics will be extracted with the CK tool, organised through an ETL pipeline in Power BI and compared with non-parametric tests (Kruskal--Wallis, Mann--Whitney, Bonferroni correction) in reproducible Python/R notebooks.
\end{abstract}

\begin{resumo}
A qualidade interna de \emph{software} orientado a objetos pode ser observada por métricas estruturais de acoplamento, complexidade, herança, coesão e tamanho. No ecossistema Java, o Spring Boot é um \emph{framework} de ampla adoção; contudo, ainda são escassas evidências longitudinais sobre como a qualidade estrutural de projetos públicos no GitHub evolui e se a Era da IA generativa (2022--atual) coincide com alterações mensuráveis nessas métricas. O objetivo deste trabalho é analisar essa evolução em quatro coortes temporais definidas pela data de criação dos repositórios. As métricas serão extraídas com a ferramenta CK, organizadas em um \emph{pipeline} de ETL no Power BI e comparadas por testes não paramétricos (Kruskal--Wallis, Mann--Whitney e correção de Bonferroni) em \emph{notebooks} reproduzíveis em Python/R.
\end{resumo}

\begin{tcolorbox}
\footnotesize
\textbf{Bacharelado em Engenharia de \emph{Software} - PUC Minas\\
Trabalho de Conclusão de Curso (TCC)} \\

\indent Orientador de conteúdo (TCC I): Danilo Maia - danilomaia@pucminas.br \\
Orientador de conteúdo (TCC I): João Pedro Batisteli - joaobatisteli@pucminas.br\\
Orientador de conteúdo (TCC I): Leonardo Vilela - leonardocardoso@pucminas.br \\
Orientador acadêmico (TCC I): Cleiton Tavares - cleitontavares@pucminas.br \\
Orientador do TCC II: (A ser definido no próximo semestre)\\ \\
Belo Horizonte, 24 de agosto de 2026.
\end{tcolorbox}

\section{Introdução}
\label{sec:introducao}

A qualidade interna de \emph{software} descreve propriedades estruturais do código-fonte que podem ser observadas sem a execução do programa, tais como modularidade, coesão, acoplamento e facilidade de manutenção \cite{iso25010,pressman:21}. Em sistemas orientados a objetos, essas propriedades têm sido operacionalizadas pela suíte de métricas proposta por Chidamber e Kemerer, que quantifica acoplamento, complexidade, herança e coesão em nível de classe \cite{chidamber:94}. No ecossistema Java, o \emph{framework} Spring Boot consolidou convenções de injeção de dependências, organização em camadas e autoconfiguração, tornando-se um recorte representativo para estudos empíricos sobre estrutura de código em repositórios públicos \cite{sun:22}. A combinação entre métricas clássicas e a ampla disponibilidade de projetos Spring Boot no GitHub permite examinar, com evidências extraídas do próprio código, como a comunidade de desenvolvimento organiza suas aplicações.

Apesar da consolidação dessas métricas e da popularidade do Spring Boot, a maior parte das investigações empíricas sobre qualidade estrutural em Java adota recortes transversais, isto é, analisa um conjunto de sistemas em um único instante, sem contrastar gerações distintas de projetos ao longo do tempo \cite{child:19,elemam:01}. Paralelamente, a partir de 2022, o desenvolvimento passou a incorporar de forma intensiva assistentes de inteligência artificial generativa, como GitHub Copilot, ChatGPT, Claude Code e Cursor, o que alterou o ambiente de produção de código \cite{peng:23,yetistiren:23}. \textbf{O problema investigado neste trabalho é a ausência de evidências empíricas longitudinais sobre como a qualidade estrutural de projetos Spring Boot hospedados no GitHub evolui ao longo de diferentes épocas históricas e se a Era da IA assistida introduziu alterações significativas nessas métricas.}

A relevância de tratar esse problema decorre de três fatores. Primeiro, a qualidade interna está associada ao esforço de manutenção e à viabilidade de longo prazo dos sistemas \cite{pressman:21,iso25010}; portanto, desconhecer a trajetória estrutural de um ecossistema tão difundido quanto o Spring Boot limita decisões de prática e de pesquisa. Segundo, a mineração de repositórios de \emph{software} no GitHub oferece base empírica em larga escala, desde que critérios de seleção, exclusão de \emph{forks} e maturidade temporal sejam observados \cite{hassan:08,kalliamvakou:16}. Terceiro, evidências recentes mostram que assistentes de IA reduzem o tempo de conclusão de tarefas e alteram a qualidade do código gerado em experimentos controlados \cite{peng:23,dakhel:23,liu:23}, mas ainda não esclarecem se esses efeitos se manifestam na estrutura de projetos reais publicados na plataforma. Sem uma análise por coortes, permanece indefinido se a Era da IA coincide com estabilidade, melhoria ou degradação das métricas estruturais.

O \textbf{objetivo geral deste trabalho é analisar longitudinalmente a evolução da qualidade estrutural de projetos Spring Boot \emph{open source} hospedados no GitHub, organizados em quatro coortes temporais, investigando o impacto da Era da IA assistida sobre métricas orientadas a objetos.} Os objetivos específicos são: (i)~caracterizar a evolução do acoplamento entre classes ao longo das quatro épocas; (ii)~caracterizar a evolução da complexidade e do tamanho de código; (iii)~caracterizar a evolução da herança e da coesão interna das classes; e (iv)~comparar estatisticamente o período imediatamente anterior à IA generativa (2018--2021) com a Era da IA assistida (2022--atual). A operacionalização desses objetivos no paradigma \emph{Goal-Question-Metric} é apresentada na Seção~\ref{sec:metodologia}.

Espera-se obter, como resultados, um retrato empírico das distribuições de acoplamento, complexidade, herança, coesão e tamanho em cada coorte, com testes não paramétricos que indiquem se as diferenças entre épocas são estatisticamente significativas, em especial no contraste entre o período pré-IA e a Era da IA assistida. Também se espera consolidar um pacote de replicação contendo \emph{scripts} de coleta, extração de métricas, notebooks de análise e artefatos de visualização, de modo a tornar o protocolo reproduzível no GitHub.

Este artigo está organizado da seguinte forma: a Seção~\ref{sec:fund} apresenta a fundamentação teórica; a Seção~\ref{sec:trabalhos} discute os trabalhos relacionados; a Seção~\ref{sec:metodologia} descreve o tipo de pesquisa, o ambiente, a amostragem, as métricas, os instrumentos, os experimentos já realizados e o cronograma do TCC~II; e o Apêndice~\ref{ap:visualizacoes} documenta o planejamento das visualizações longitudinais.

\section{Fundamentação Teórica}
\label{sec:fund}

Esta seção apresenta os conceitos consolidados que embasam o estudo: qualidade interna de \emph{software}, a suíte de métricas de Chidamber e Kemerer, o paradigma \emph{Goal-Question-Metric}, a mineração de repositórios e a inteligência artificial generativa no recorte da Era da IA assistida.

\subsection{Qualidade interna e métricas orientadas a objetos}
\label{subsec:fund-qualidade}

A norma \emph{ISO/IEC 25010} define qualidade de \emph{software} como um conjunto de características que expressam o grau em que um sistema atende a necessidades explícitas e implícitas \cite{iso25010}. Nesse modelo, a qualidade interna refere-se a propriedades do código-fonte observáveis por análise estática, incluindo modularidade, coesão e facilidade de manutenção. Pressman e Maxim destacam que sistemas com melhor qualidade interna tendem a exigir menor esforço de evolução e a permanecer viáveis por mais tempo \cite{pressman:21}.

Para tornar essas características mensuráveis em código orientado a objetos, Chidamber e Kemerer propuseram um conjunto clássico de indicadores em nível de classe \cite{chidamber:94}. Acoplamento entre objetos (\emph{Coupling Between Objects} --- CBO) descreve o grau de dependência de uma classe em relação a outras. Métodos ponderados por classe (\emph{Weighted Methods per Class} --- WMC) agrega a complexidade dos métodos. Profundidade da árvore de herança (\emph{Depth of Inheritance Tree} --- DIT) e número de filhos (\emph{Number of Children} --- NOC) caracterizam a hierarquia. Resposta de uma classe (\emph{Response for a Class} --- RFC) estima o conjunto de métodos potencialmente acionados a partir de uma mensagem. Falta de coesão em métodos (\emph{Lack of Cohesion in Methods} --- LCOM) indica o quanto os métodos de uma classe operam sobre conjuntos disjuntos de atributos. Essas métricas passaram a ser amplamente empregadas para caracterizar projeto, manutenção e falhas em sistemas orientados a objetos \cite{singh:12,radjenovic:13,aniche:16}. A forma operacional em que serão extraídas neste trabalho, incluindo a ferramenta e as variantes adotadas, é especificada na Seção~\ref{sec:metodologia}.

\subsection{Goal-Question-Metric}
\label{subsec:fund-gqm}

O paradigma \emph{Goal-Question-Metric} (GQM) organiza avaliações de \emph{software} em três níveis: o objetivo do estudo (\emph{goal}), as perguntas que o detalham (\emph{questions}) e as métricas que as respondem (\emph{metrics}) \cite{basili:94}. A abordagem evita coletar indicadores sem vínculo com a intenção da pesquisa: primeiro define-se o que se deseja aprender; em seguida, traduz-se esse objetivo em perguntas mensuráveis; por fim, escolhem-se as medidas adequadas. Neste trabalho, o GQM liga o objetivo de analisar a evolução da qualidade estrutural em projetos Spring Boot às dimensões de acoplamento, complexidade, herança, coesão e tamanho, operacionalizadas na Seção~\ref{sec:metodologia}.

\subsection{Mineração de Repositórios de Software}
\label{subsec:fund-msr}

A Mineração de Repositórios de \emph{Software} (\emph{Mining Software Repositories} --- MSR) analisa artefatos reais --- código-fonte, histórico de versões e metadados de plataformas --- para identificar padrões de desenvolvimento, evolução e qualidade \cite{hassan:08}. Em vez de apoiar-se apenas em relatos ou documentação, a MSR observa evidências extraídas dos próprios repositórios. No GitHub, essa estratégia exige filtros de relevância, exclusão de \emph{forks} e consideração de um intervalo de maturação após a criação do repositório, a fim de reduzir ruído e projetos ainda imaturos \cite{kalliamvakou:14,kalliamvakou:16}. Esses cuidados conceituais orientam a amostragem descrita na metodologia.

\subsection{Inteligência artificial generativa}
\label{subsec:fund-ia}

Neste trabalho, inteligência artificial generativa designa sistemas baseados em modelos de linguagem de grande escala, treinados em corpora de texto e código, capazes de produzir ou completar trechos de programa a partir de contexto e de \emph{prompts}. Ferramentas como GitHub Copilot, ChatGPT, Claude Code e Cursor ilustram esse paradigma e delimitam o que se denomina Era da IA assistida, adotada como quarta época histórica do estudo. Experimentos controlados indicam ganhos de produtividade e variações na qualidade do código sugerido por esses assistentes \cite{peng:23,dakhel:23}, o que motiva investigar se a estrutura interna de projetos Spring Boot publicados no GitHub apresenta sinais observáveis nesse período.

\section{Trabalhos Relacionados}
\label{sec:trabalhos}

Esta seção discute cinco estudos empíricos da engenharia de \emph{software} que tratam da qualidade de código assistido por IA, da corretude de artefatos gerados por modelos de linguagem ou da interpretação de métricas de acoplamento em Java. Para cada trabalho, apresentam-se o objetivo e o método, os resultados pertinentes a este TCC e a relação com a proposta --- complementar, comparar ou reutilizar o mesmo tipo de evidência. Os três primeiros artigos foram publicados em 2023 e estão diretamente ligados ao recorte da Era da IA assistida.

A qualidade do código produzido por assistentes foi avaliada de forma experimental com o conjunto HumanEval, comparando GitHub Copilot, Amazon CodeWhisperer e ChatGPT quanto a validade, correção, segurança, confiabilidade e manutenibilidade \cite{yetistiren:23}. O método consistiu em gerar soluções a partir de enunciados e inspecionar o código resultante com métricas de qualidade e critérios de aceitação dos testes do \emph{benchmark}. Os resultados indicam taxas de correção de 65,2\% para o ChatGPT, 46,3\% para o Copilot e 31,1\% para o CodeWhisperer, além de dívida técnica média associada a \emph{code smells} da ordem de minutos por trecho. Esse estudo é complementar ao presente trabalho: enquanto avalia trechos gerados em tarefas isoladas, esta pesquisa observa a estrutura de repositórios Spring Boot reais ao longo do tempo, permitindo contrastar o período pré-IA com a Era da IA assistida.

A capacidade do GitHub Copilot como programador-par foi examinada em dois conjuntos de tarefas: problemas algorítmicos fundamentais e um corpus de soluções humanas, corretas e defeituosas \cite{dakhel:23}. O método combinou geração repetida de soluções, verificação de funcionalidade e comparação com implementações de programadores. Os resultados mostram que o Copilot produz soluções para a maioria dos problemas, porém parte delas é defeituosa ou não reproduzível; a taxa de correção humana superou a do assistente, embora o código gerado com falhas tenda a exigir menor esforço de reparo. Os autores concluem que a ferramenta pode ser um ativo para desenvolvedores experientes e um passivo para novatos. O presente trabalho não replica o desenho de tarefas algorítmicas; em vez disso, complementa essa evidência ao perguntar se a adoção massiva de assistentes, a partir de 2022, coincide com mudanças nas métricas estruturais de projetos Java/Spring Boot publicados no GitHub.

A correção funcional de código gerado por modelos de linguagem, incluindo o ChatGPT, foi avaliada com o \emph{framework} EvalPlus, que amplia o HumanEval com uma quantidade substancialmente maior de casos de teste obtidos por geração automática e mutação \cite{liu:23}. O método consistiu em reexecutar 26 modelos populares sobre HumanEval e HumanEval+, medindo \emph{pass@k}. Os resultados mostram quedas de até 19,3\%--28,9\% em \emph{pass@k} quando os testes adicionais são considerados, e alteração no ranqueamento de modelos --- inclusive casos em que o ChatGPT deixa de superar alternativas de código aberto. A relação com este trabalho é de complementaridade: os resultados de \cite{liu:23} demonstram que a qualidade funcional do código gerado por IA é frequentemente superestimada em \emph{benchmarks} curtos; o presente estudo desloca o foco da correção de funções sintéticas para a qualidade estrutural de sistemas Spring Boot completos, em coortes históricas.

O impacto da injeção de dependências sobre a manutenibilidade em projetos Java foi investigado com duas métricas ponderadas --- DCE e DCBO --- implementadas no CKJM-Analyzer, extensão da ferramenta CKJM \cite{sun:22}. O método aplicou análise estática a um conjunto de projetos \emph{open source}, distinguindo acoplamento ``suave'' (via injeção) de acoplamento ``rígido'' (instanciação direta). Os resultados indicam que a injeção de dependências reduz o acoplamento efetivo quando este é ponderado, o que métricas clássicas de CBO tendem a não capturar integralmente. O presente trabalho compartilha o objeto de estudo --- código Java com forte presença de injeção, típica do Spring Boot ---, porém não propõe métricas novas: utiliza a suíte CK de forma consistente entre épocas para observar evolução estrutural, reconhecendo a limitação apontada por \cite{sun:22} na interpretação do acoplamento.

O efeito do GitHub Copilot sobre a produtividade de desenvolvedores profissionais foi medido em experimento controlado, com tarefa de implementação cronometrada e grupo de controle sem o assistente \cite{peng:23}. O método alocou participantes de forma aleatória e comparou tempo de conclusão e taxa de término. Os resultados apontam redução substancial do tempo necessário para completar a tarefa (cerca de 55\% mais rápido no grupo com Copilot) e maior taxa de conclusão. Essa evidência fundamenta a delimitação da quarta época histórica (a partir de 2022) como Era da IA assistida. A diferença em relação a este trabalho é o construto medido: o experimento de \cite{peng:23} avalia desempenho individual em uma tarefa; aqui se avaliam métricas estruturais de repositórios públicos, sem reivindicar causalidade direta entre o uso de um assistente específico e a qualidade de cada projeto.

Em conjunto, os estudos revisados mostram que assistentes de IA alteram produtividade e qualidade de trechos gerados, e que métricas clássicas de acoplamento exigem cautela em sistemas com injeção de dependências. Nenhum deles, contudo, descreve a trajetória longitudinal da qualidade estrutural de projetos Spring Boot no GitHub. Essa lacuna é o ponto de partida da metodologia apresentada a seguir.

\section{Materiais e Métodos}
\label{sec:metodologia}

Esta seção descreve o tipo de pesquisa, o ambiente e os instrumentos, a amostragem, as métricas de avaliação, os procedimentos estatísticos, os experimentos já realizados e o cronograma do TCC~II. O desenho é classificado como \textbf{empírico}, porque se baseia em repositórios reais do GitHub; \textbf{quantitativo}, porque emprega métricas estruturais e testes estatísticos formais; e \textbf{longitudinal}, porque organiza a amostra em quatro coortes temporais definidas pela data de criação do repositório, permitindo contrastar gerações distintas de projetos Spring Boot e isolar o salto entre o período pré-IA (E3) e a Era da IA assistida (E4). A classificação justifica-se pelo objetivo de obter evidências observáveis, comparáveis entre épocas, e não pelo acompanhamento do mesmo repositório ao longo de toda a sua vida.

Importa reconhecer uma limitação do desenho por coortes: repositórios de épocas mais antigas tiveram mais tempo de evolução orgânica do que os criados recentemente. Por isso, diferenças entre épocas não serão interpretadas automaticamente como melhoria ou piora de qualidade, mas como variações associadas à coorte, com controle explícito do tamanho do projeto (LOC) e da versão do Spring Boot.

A Figura~\ref{fig:metodologia} resume o fluxo em três etapas: coleta dos repositórios; extração das métricas com a ferramenta CK; e consolidação para análise estatística e visualização.

\begin{figure}[H]
    \centering
    \begin{tikzpicture}[
        font=\small,
        box/.style={
            draw,
            rounded corners=3pt,
            align=center,
            inner sep=7pt,
            text width=0.27\linewidth,
            minimum height=2.6cm,
            thick
        },
        arr/.style={-{Latex}, thick, shorten >=2pt, shorten <=2pt}
    ]
        \node[box] (coleta) {\textbf{1. Coleta}\\GitHub GraphQL\\filtros Java/Spring Boot\\coortes E1--E4};
        \node[box, right=0.45cm of coleta] (extracao) {\textbf{2. Extração}\\Ferramenta CK\\CBO, RFC, WMC,\\DIT, NOC, LCOM\_HS, LOC};
        \node[box, right=0.45cm of extracao] (analise) {\textbf{3. Análise}\\ETL no Power BI\\testes estatísticos\\visualização};
        \draw[arr] (coleta) -- (extracao);
        \draw[arr] (extracao) -- (analise);
    \end{tikzpicture}
    \caption{Visão geral da metodologia: coleta, extração de métricas e análise.}
    \label{fig:metodologia}
    {\footnotesize Fonte: Elaborada pelos autores.}
\end{figure}

\subsection{Ambiente e instrumentos}
\label{subsec:ambiente}

A coleta será realizada em Python, com a \emph{GitHub GraphQL API}, paginação e controle de \emph{rate limit}. Essa interface permite recuperar metadados específicos (\texttt{createdAt}, estrelas, linguagens, \emph{fork}) em uma única consulta estruturada, o que torna a amostragem mais reprodutível do que uma sequência de chamadas REST. A análise estrutural será executada com a ferramenta CK, versão 0.7.1, sobre JVM em versão fixa, documentada no pacote de replicação \cite{aniche:ck}. A CK foi escolhida por calcular, de forma automatizada, a suíte de Chidamber e Kemerer em código Java e por ser amplamente utilizada em estudos de MSR. O Power BI será o ambiente de ETL e de \emph{dashboards} exploratórios. A inferência estatística e as figuras finais reproduzíveis serão conduzidas em \emph{notebooks} Python (\texttt{scipy}, \texttt{statsmodels}) ou R. Diferentes ferramentas de métricas podem implementar o mesmo indicador de modos distintos, em especial o acoplamento \cite{child:19}; por isso a versão da CK, da JVM e os parâmetros de execução permanecerão constantes em todas as épocas.

\subsection{Amostragem e procedimentos}
\label{subsec:amostragem}

A amostra será formada por repositórios públicos escritos principalmente em Java e com uso explícito de Spring Boot (dependências \texttt{spring-boot-starter} em \texttt{pom.xml} ou \texttt{build.gradle}). Serão considerados apenas projetos não forkeados, com no mínimo 10 estrelas, 2.000 linhas de código efetivas e 10 arquivos Java. Esses critérios reduzem tutoriais triviais e repositórios sem base de código adequada para análise estrutural. Será aplicado um \emph{delay} de mineração de 12 meses \cite{kalliamvakou:16}, para evitar que repositórios muito recentes sejam tratados como representativos de uma época ainda em consolidação. Pretende-se coletar 500 repositórios por época, totalizando 2.000 projetos. Os metadados incluem identificador, URL, estrelas, \texttt{createdAt}, \texttt{updatedAt}, estimativa de tamanho, linguagens e a versão do Spring Boot declarada no arquivo de \emph{build}.

Cada repositório será associado a exatamente uma época com base em \texttt{createdAt} (UTC), conforme a Tabela~\ref{tab:epocas}. Após a seleção, os repositórios elegíveis serão clonados e analisados estaticamente pela CK. Os resultados serão armazenados em CSV com identificação do repositório e da época. Caso ocorram falhas de clonagem ou de extração, o projeto será substituído por outro da mesma época que respeite os mesmos critérios. Os CSV serão consolidados em um processo de ETL no Power BI, restrito a limpeza, modelagem dimensional e visualização exploratória. As agregações por repositório serão exportadas para CSV/Parquet que alimentarão a análise estatística formal. O objetivo não é acompanhar a evolução completa de cada repositório, e sim construir um retrato da coorte a partir de projetos criados na janela histórica correspondente.

\begin{table}[H]
    \centering
    \caption{Épocas históricas: delimitação operacional das coortes no GitHub.}
    \label{tab:epocas}
    \footnotesize
    \renewcommand{\arraystretch}{1.2}
    \begin{tabularx}{\linewidth}{|c|c|c|>{\raggedright\arraybackslash}p{0.18\linewidth}|>{\raggedright\arraybackslash}X|}
        \hline
        \textbf{Época} & \textbf{Início} & \textbf{Fim} & \textbf{Denominação} & \textbf{Marcos representativos} \\
        \hline
        E1 & 01/01/2009 & 31/12/2013 & Colaboração emergente &
        Consolidação inicial do GitHub e do ecossistema \emph{open source} Java. \\
        \hline
        E2 & 01/01/2014 & 31/12/2017 & Estruturação e \emph{frameworks} &
        Popularização do Spring Boot; difusão de Docker e CI. \\
        \hline
        E3 & 01/01/2018 & 31/12/2021 & Maturidade colaborativa &
        Consolidação de práticas de qualidade e revisão por pares. \\
        \hline
        E4 & 01/01/2022 & data da coleta & Era da IA assistida &
        GitHub Copilot, Claude Code, Cursor, ChatGPT e adoção de IA generativa. \\
        \hline
    \end{tabularx}
    \vspace{4pt}
    {\footnotesize Fonte: Elaborada pelos autores. Critério de alocação: campo \texttt{createdAt} do repositório no GitHub (limites inclusivos).}
\end{table}

\subsection{Métricas de avaliação}
\label{subsec:metricas-avaliacao}

As métricas extraídas pela CK, em nível de classe e de método, são CBO, RFC, WMC, DIT, NOC, LCOM\_HS e LOC. Para LCOM, utiliza-se explicitamente a variante de Henderson-Sellers (LCOM\_HS), exportada na coluna \texttt{lcom} pela CK, distinta de LCOM1, LCOM2 e LCOM3 \cite{child:19,aniche:ck}. Após a sumarização, cada repositório será representado por médias: \texttt{avg\_cbo}, \texttt{avg\_rfc}, \texttt{avg\_wmc}, \texttt{avg\_loc\_method}, \texttt{avg\_dit}, \texttt{avg\_noc} e \texttt{avg\_lcom}. CBO e RFC observam acoplamento; WMC e LOC, complexidade e tamanho; DIT e NOC, herança; LCOM\_HS, coesão. A escolha justifica-se por duas razões: são indicadores clássicos e comparáveis entre projetos \cite{chidamber:94} e mapeiam diretamente os objetivos específicos da introdução. Métricas de complexidade e acoplamento são sensíveis ao tamanho das classes e dos métodos \cite{elemam:01}; por isso o LOC entra tanto como dimensão de análise quanto como variável de controle.

A Tabela~\ref{tab:gqm-objetivos} unifica objetivos específicos, questões de pesquisa e métricas, conforme o GQM (Seção~\ref{subsec:fund-gqm}).

\begin{table}[H]
    \centering
    \caption{Objetivos específicos, questões de pesquisa e métricas (GQM).}
    \label{tab:gqm-objetivos}
    \footnotesize
    \renewcommand{\arraystretch}{1.2}
    \begin{tabularx}{\linewidth}{|>{\raggedright\arraybackslash}p{0.26\linewidth}|c|>{\raggedright\arraybackslash}p{0.14\linewidth}|>{\raggedright\arraybackslash}X|>{\raggedright\arraybackslash}p{0.14\linewidth}|}
        \hline
        \textbf{Objetivo específico} & \textbf{Q} & \textbf{Dimensão} & \textbf{Questão de pesquisa} & \textbf{Métricas} \\
        \hline
        Caracterizar acoplamento entre épocas &
        Q1 & Acoplamento &
        Como a dependência entre classes evoluiu ao longo das quatro
        épocas? & CBO, RFC \\
        \hline
        Caracterizar complexidade e tamanho entre épocas &
        Q2 & Complexidade &
        A densidade lógica dos métodos e o tamanho do código se alteraram
        com o tempo e na Era da IA? & WMC, LOC \\
        \hline
        Caracterizar herança e coesão entre épocas &
        Q3 & Herança e coesão &
        As hierarquias permaneceram rasas e a coesão interna das classes
        se manteve estável? & DIT, NOC, LCOM\_HS \\
        \hline
        Contrastar E3 e E4 (Era da IA) &
        Q4 & Contraste IA &
        O período pós-2022 difere de 2018--2021 nas métricas
        estruturais? & Todas \\
        \hline
    \end{tabularx}
    \vspace{4pt}
    {\footnotesize Fonte: Elaborada pelos autores.}
\end{table}

\subsection{Métodos estatísticos}
\label{subsec:estatistica}

Além do filtro mínimo de 2.000 LOC, serão reportadas as distribuições de LOC e de versão do Spring Boot por época. Antes das comparações principais, aplicar-se-á o teste de Kruskal--Wallis sobre LOC entre épocas. Se diferenças forem detectadas, as análises estruturais serão complementadas por estratificação por faixas de LOC e por análise de sensibilidade com métricas normalizadas (razões por classe ou por KLOC). A versão do Spring Boot será registrada por repositório e tratada como variável confundidora na discussão.

Para cada métrica agregada, aplicar-se-á o teste de Kruskal--Wallis às quatro épocas. Havendo significância global, seguir-se-á o teste de Dunn com correção de Bonferroni para comparações par a par. Para a hipótese da Era da IA, empregar-se-á o teste de Mann--Whitney entre E3 e E4. A correção para comparações múltiplas adotará $\alpha_{\mathrm{adj}} = \alpha / m$, com $\alpha = 0{,}05$ e $m$ igual ao número de testes da família: $m = 7$ no contraste E3 versus E4; $m = 42$ no conjunto completo de pares de épocas após Kruskal--Wallis (seis pares vezes sete métricas). Serão reportados \emph{p-valores} ajustados, tamanho de efeito (\emph{rank-biserial} ou $\delta$ de Cliff) e intervalos descritivos (mediana e amplitude interquartil). As visualizações previstas constam do Apêndice~\ref{ap:visualizacoes}.

\subsection{Experimentos já realizados}
\label{subsec:piloto}

Um estudo piloto já foi executado para validar o \emph{pipeline} de extração. \emph{Scripts} em Python clonam repositórios Java públicos, invocam a CK 0.7.1 e consolidam CBO, RFC, WMC, DIT, NOC, LOC e indicadores auxiliares em CSV. O artefato \texttt{analise\_repositorios\_java.csv} reúne a análise de 520 repositórios, realizada em novembro de 2025, incluindo projetos Spring Boot de ampla adoção, como \texttt{spring-projects/spring-boot} (média de CBO 5,45; WMC 4,96; DIT 1,29) e \texttt{xkcoding/spring-boot-demo} (CBO 5,05; WMC 2,77; DIT 1,13). O piloto confirma que o ambiente (Python, Git, JVM e CK) está operacional e que as métricas previstas são extraíveis em escala. Não substitui a amostra final: os repositórios ainda não foram filtrados por dependência Spring Boot nem alocados às quatro épocas históricas. Esses recortes serão aplicados na coleta GraphQL descrita acima.

\subsection{Cronograma do TCC II}
\label{subsec:cronograma}

O trabalho é desenvolvido em dupla. As Tabelas~\ref{tab:cronograma-gabriel} e~\ref{tab:cronograma-guilherme} apresentam cronogramas quinzenais independentes e complementares, da primeira quinzena de agosto de 2026 à segunda quinzena de novembro de 2026. Gabriel Ferreira concentra-se na amostragem, na mineração GraphQL e no ETL; Guilherme Augusto concentra-se na extração CK, na inferência estatística e nas figuras finais.

\begin{table}[H]
\centering
\caption{Cronograma quinzenal --- Gabriel Ferreira.}
\label{tab:cronograma-gabriel}
\footnotesize
\renewcommand{\arraystretch}{1.15}
\begin{tabularx}{\linewidth}{|>{\raggedright\arraybackslash}X|c|c|c|c|c|c|c|c|}
\hline
\textbf{Tarefa} & \textbf{A1} & \textbf{A2} & \textbf{S1} & \textbf{S2} & \textbf{O1} & \textbf{O2} & \textbf{N1} & \textbf{N2} \\
\hline
Critérios de amostragem e consulta GraphQL & X & X & & & & & & \\
\hline
Mineração E1--E2 e aplicação dos filtros & & & X & & & & & \\
\hline
Mineração E3--E4 e versão do Spring Boot & & & & X & & & & \\
\hline
ETL no Power BI e limpeza dos dados & & & & & X & X & & \\
\hline
Dashboards exploratórios e exportação CSV & & & & & & X & X & \\
\hline
Documentação da coleta e pacote de replicação & & & & & & & X & X \\
\hline
\end{tabularx}
\vspace{4pt}
{\footnotesize Fonte: Elaborada pelos autores. A1--A2: agosto; S1--S2: setembro; O1--O2: outubro; N1--N2: novembro de 2026.}
\end{table}

\begin{table}[H]
\centering
\caption{Cronograma quinzenal --- Guilherme Augusto.}
\label{tab:cronograma-guilherme}
\footnotesize
\renewcommand{\arraystretch}{1.15}
\begin{tabularx}{\linewidth}{|>{\raggedright\arraybackslash}X|c|c|c|c|c|c|c|c|}
\hline
\textbf{Tarefa} & \textbf{A1} & \textbf{A2} & \textbf{S1} & \textbf{S2} & \textbf{O1} & \textbf{O2} & \textbf{N1} & \textbf{N2} \\
\hline
Ambiente CK/JVM e validação do piloto & X & & & & & & & \\
\hline
\textit{Pipeline} de clonagem e extração CK & & X & X & & & & & \\
\hline
Extração CK e agregação \texttt{avg\_*} (E1--E4) & & & X & X & & & & \\
\hline
Testes Kruskal--Wallis, Dunn e Mann--Whitney & & & & & X & X & & \\
\hline
Figuras finais e interpretação estatística & & & & & & X & X & \\
\hline
Pacote de replicação da análise e revisão do texto & & & & & & & X & X \\
\hline
\end{tabularx}
\vspace{4pt}
{\footnotesize Fonte: Elaborada pelos autores. A1--A2: agosto; S1--S2: setembro; O1--O2: outubro; N1--N2: novembro de 2026.}
\end{table}


\bibliographystyle{sbc}
\bibliography{sbc-template}

\appendix
\section*{Apêndice A --- Planejamento das visualizações longitudinais}
\label{ap:visualizacoes}

As Tabelas~\ref{tab:planejamento-figuras} e~\ref{tab:planejamento-figuras-cont}
apresentam o planejamento das visualizações que serão geradas durante a
execução do TCC~II, relacionando cada figura prevista à respectiva questão
de pesquisa, à métrica principal e ao tipo de visualização adotado. O
conjunto foi desenhado para permitir a comparação entre as quatro épocas
históricas E1 (2009--2013), E2 (2014--2017), E3 (2018--2021) e E4
(2022--atual), conforme delimitação operacional na
Seção~\ref{sec:metodologia}, e para evidenciar \emph{(i)}~tendências
evolutivas de longo prazo (E1~$\rightarrow$~E3) e \emph{(ii)}~o eventual
impacto da Era da IA assistida (E3~$\rightarrow$~E4).

\begin{table}[htbp]
\centering
\caption{Planejamento das visualizações por questão de pesquisa (parte~1).}
\label{tab:planejamento-figuras}
\scriptsize
\setlength{\tabcolsep}{3pt}
\renewcommand{\arraystretch}{1.05}
\begin{tabularx}{\linewidth}{|c|c|>{\raggedright\arraybackslash}p{0.20\linewidth}|>{\raggedright\arraybackslash}p{0.18\linewidth}|>{\raggedright\arraybackslash}X|}
\hline
\textbf{Fig.} & \textbf{Q} & \textbf{Métrica} & \textbf{Visualização} & \textbf{Descrição} \\
\hline
F1 & Q1 & CBO &
\emph{Boxplot} &
Acoplamento médio entre classes por época (mediana, dispersão e
\emph{outliers}). \\
\hline
F2 & Q1 & RFC &
\emph{Boxplot} &
Métodos potencialmente acionados por classe em cada época. \\
\hline
F3 & Q1 & CBO, RFC &
Série temporal &
Medianas de CBO e RFC ao longo das quatro épocas. \\
\hline
F4 & Q2 & WMC &
\emph{Boxplot} &
Complexidade lógica agregada por classe em cada época. \\
\hline
F5 & Q2 & \texttt{avg\_loc\_method} &
\emph{Boxplot} &
Tamanho médio dos métodos; hipótese de métodos maiores na Era da IA. \\
\hline
\end{tabularx}
\end{table}

\begin{table}[htbp]
\centering
\caption{Planejamento das visualizações por questão de pesquisa (parte~2).}
\label{tab:planejamento-figuras-cont}
\scriptsize
\setlength{\tabcolsep}{3pt}
\renewcommand{\arraystretch}{1.05}
\begin{tabularx}{\linewidth}{|c|c|>{\raggedright\arraybackslash}p{0.20\linewidth}|>{\raggedright\arraybackslash}p{0.18\linewidth}|>{\raggedright\arraybackslash}X|}
\hline
\textbf{Fig.} & \textbf{Q} & \textbf{Métrica} & \textbf{Visualização} & \textbf{Descrição} \\
\hline
F6 & Q2 & \texttt{avg\_wmc}, \texttt{avg\_loc\_method} &
Dispersão &
Correlação entre tamanho e complexidade médios, colorida por época. \\
\hline
F7 & Q3 & DIT &
\emph{Violin plot} &
Profundidade média da hierarquia de herança por época. \\
\hline
F8 & Q3 & NOC &
\emph{Violin plot} &
Número médio de subclasses diretas por classe. \\
\hline
F9 & Q3 & LCOM\_HS &
\emph{Boxplot} &
Ausência média de coesão por repositório em cada época. \\
\hline
F10 & Q4 & Todas &
Painel comparativo E3 vs.\ E4 &
Contraste da Era da IA, com \emph{p-valores} ajustados (Mann--Whitney). \\
\hline
\end{tabularx}
\end{table}

{\footnotesize Elaborado pelos autores. Os artefatos finais (PNG,
\emph{notebooks} Python/R e arquivo \texttt{.pbix} do Power~BI) serão
publicados no pacote de replicação previsto para o TCC~II.}

\medskip
Para garantir a comparabilidade entre figuras, todas as visualizações
adotarão a mesma paleta de cores por época, escala consistente
entre painéis e identificação explícita das épocas E1 (2009--2013),
E2 (2014--2017), E3 (2018--2021) e E4 (2022--atual). Sempre que
cabível, os \emph{p-valores} ajustados (Bonferroni) dos testes de
Kruskal--Wallis e Mann--Whitney (E3 \emph{vs.} E4) serão reportados no
próprio painel ou em legenda associada.


\end{document}